\subsection{Harness Implementation Architecture}
\label{app:harness_architecture}
\label{app:harness_implementation}  % Legacy label for backward compatibility

This appendix provides implementation details of the Python-based benchmarking harness that orchestrated all experimental runs reported in \cref{sec:results}. The harness is organized into modular components with clear separation of concerns, enabling independent development, testing, and replication. All code is maintained in a single repository with versioned dependencies to ensure reproducible builds (\cref{sec:reproducibility_guide}).

\subsubsection{System Architecture and Components}

The harness comprises six cooperating subsystems: configuration management, execution scheduling, provider adapters, telemetry and cost accounting, persistence and export, and web-based monitoring.

\paragraph{Configuration management.} Model definitions reside in \texttt{models.json}, a validated JSON file specifying per-model metadata including canonical model identifier (\texttt{model}), version tag (\texttt{model\_version}), category labels (\texttt{categories}), rate limit parameters (\texttt{rate\_limit\_per\_min}, \texttt{cooldown\_s}), default temperature (\texttt{temperature}), and enabled status (\texttt{enabled}). System prompt templates are stored in \texttt{system\_prompts.json} with Jinja2 variables for dynamic instantiation. User prompts are organized in \texttt{test\_prompts.json}, with fields \texttt{prompt\_id}, \texttt{category}, \texttt{prompt\_text}, optional \texttt{teacher\_context}, and \texttt{expected\_keywords} for automated fidelity checks (\cref{subsec:automated_metrics}). A JSON Schema validator runs at startup to detect malformed configurations before execution begins.

Sensitive credentials are externalized via environment variables loaded from a \texttt{.env} file excluded from version control per \cref{subsubsec:security_privacy}. The loader checks for required keys (\texttt{OPENAI\_API\_KEY}, \texttt{ANTHROPIC\_API\_KEY}, etc.) and emits actionable diagnostics when keys are missing or appear malformed.

\paragraph{Execution scheduling.} The scheduler maintains a global worker pool (default: 4 workers) and per-provider cooldown timers to enforce rate limits described in \cref{subsubsec:core_functionality}. Before dispatching a call, the scheduler verifies that (i) the target provider's cooldown has elapsed, (ii) the global concurrency quota is not exhausted, and (iii) the model's adapter is available. If any constraint is violated, the task is deferred and retried after a configurable backoff interval. This design ensures symmetric treatment across providers, mitigating bias from transient failures (\cref{subsec:threats_limitations}).

The scheduler implements a fail-fast policy for persistent configuration errors (e.g., unsupported model IDs), logging a standardized \texttt{error\_code} and halting execution to prevent incomplete or biased runs. Transient errors (HTTP 429, network timeouts) trigger exponential backoff with jitter (base delay 2s, multiplier 2.0, max 5 retries) as documented in \cref{subsubsec:core_functionality}.

\paragraph{Provider adapters.} Each LLM provider is supported by a dedicated adapter module that implements the standardized interface contract specified in \cref{subsubsec:adapter_interface}. Adapters encapsulate provider-specific authentication, request serialization, response parsing, token counting, and error handling. All adapters return a uniform dictionary with keys \texttt{response\_text}, \texttt{tokens\_in}, \texttt{tokens\_out}, \texttt{latency\_ms}, \texttt{cost\_usd}, \texttt{http\_status}, \texttt{model\_version}, and optional \texttt{error\_code}/\texttt{error\_message}.

\subparagraph{OpenAI adapter (\texttt{openai\_adapter.py}).} Uses the official \texttt{openai} Python client (v1.14.3). Authenticates via \texttt{OPENAI\_API\_KEY}. Constructs messages as a list with \texttt{system} and \texttt{user} roles. Latency is timed via \texttt{time.perf\_counter()} bracketing the synchronous \texttt{chat.completions.create()} call. Token counts are extracted from \texttt{response.usage.prompt\_tokens} and \texttt{response.usage.completion\_tokens}. Cost is computed using posted GPT-4o Mini pricing (\$0.15/1M input tokens, \$0.60/1M output tokens) per \cref{eq:per_call_cost}. Errors are caught in a try/except block and mapped to standardized codes.

\subparagraph{Anthropic adapter (\texttt{anthropic\_adapter.py}).} Uses the \texttt{anthropic} client (v0.21.3). Authenticates via \texttt{ANTHROPIC\_API\_KEY}. System prompt is passed to the \texttt{system} parameter; user prompt to \texttt{messages} with role \texttt{user}. Token usage is read from \texttt{response.usage.input\_tokens} and \texttt{response.usage.output\_tokens}. Pricing for Claude 3 Sonnet (\$3.00/1M input, \$15.00/1M output) is hardcoded in the adapter and logged with run metadata for traceability (\cref{subsec:cost}). Model version is retrieved from \texttt{response.model}.

\subparagraph{Google adapter (\texttt{google\_adapter.py}).} Uses \texttt{google-generativeai} SDK (v0.3.2). Authenticates via \texttt{GOOGLE\_API\_KEY}. System instructions are set via \texttt{genai.GenerativeModel(system\_instruction=...)}. User prompt is passed to \texttt{generate\_content()}. Token counts are extracted from \texttt{response.usage\_metadata.prompt\_token\_count} and \texttt{response.usage\_metadata.candidates\_token\_count}. Gemini 2.5 Flash pricing (\$0.075/1M input, \$0.30/1M output) is applied. Because the SDK does not directly expose HTTP status codes for successful calls, we infer status 200 on non-exception returns and catch provider exceptions to populate \texttt{error\_code}.

\subparagraph{Cohere adapter (\texttt{cohere\_adapter.py}).} Illustrated in \cref{lst:cohere_adapter}. Uses \texttt{cohere} client (v4.38). System prompt is passed as \texttt{preamble}; user prompt as \texttt{message}. Token counts are read from \texttt{response.meta.billed\_units.input\_tokens} and \texttt{response.meta.billed\_units.output\_tokens}. Pricing for Command R 08-2024 (\$0.50/1M input, \$1.50/1M output) is hardcoded.

\subparagraph{Llama 4 Maverick adapter (\texttt{llama\_adapter.py}).} Uses Hugging Face Inference API via \texttt{huggingface\_hub.InferenceClient} (v0.20.1). Authenticates via \texttt{HUGGINGFACE\_API\_KEY}. System prompt and user prompt are concatenated with formatting tokens (\texttt{<|system|>}, \texttt{<|user|>}, \texttt{<|assistant|>}). Latency and token counts are estimated; Hugging Face does not bill per token but per inference-second, so \texttt{cost\_usd} is set to zero in the adapter and treated separately in \cref{sec:cost_modelling}. Model version is hardcoded to \texttt{meta-llama/Llama-4-Maverick-17B-128E-Instruct}.

\paragraph{Telemetry and cost accounting.}\label{tab:model_pricing} Each call's latency is measured in milliseconds using high-resolution timers. Token counts are obtained from provider response objects when available; if missing, we estimate via a fallback tokenizer aligned to the provider's documented behavior (\texttt{tiktoken} for OpenAI-compatible endpoints, \texttt{anthropic.count\_tokens()} for Anthropic, character-based approximations with tuned multipliers for others). Per-call cost is computed as specified in \cref{eq:per_call_cost}:
\begin{equation*}
\text{cost\_usd} = \frac{\text{tokens\_in}}{1{,}000{,}000} \cdot p_{\text{in}} + \frac{\text{tokens\_out}}{1{,}000{,}000} \cdot p_{\text{out}},
\end{equation*}
with pricing parameters ($p_{\text{in}}$, $p_{\text{out}}$) recorded in run metadata. Model pricing details are documented throughout this section.

\paragraph{Persistence and export.} Results are written to two destinations: (i) timestamped CSV files in \texttt{results/raw\_output/} with schema documented in \cref{app:csv_schema}, and (ii) a SQLite database (\texttt{results/benchmark.db}) with normalized tables (\texttt{runs}, \texttt{prompts}, \texttt{responses}) per \cref{app:sqlite_schema}. CSV headers and SQLite column types are validated against schema definitions at startup. Writes are buffered and flushed periodically (default every 10 records or 30s) to balance durability and performance. On graceful shutdown or fatal errors, pending buffers are flushed to ensure no data loss.

\paragraph{Web-based monitoring.} A lightweight Flask server (\texttt{dashboard.py}) provides real-time visibility during runs. Endpoints serve: (i) a run index with filters by model, category, and date; (ii) interactive latency and cost distribution plots using Plotly.js; (iii) success/failure heatmaps categorized by error class; (iv) side-by-side response viewers for qualitative inspection. The dashboard queries the SQLite database and does not modify state, allowing concurrent access during benchmark execution. This interface facilitated exploratory analysis and error triage prior to formal statistical reporting in \cref{sec:results}.

\subsubsection{Main Execution Loop}

The harness entry point (\texttt{main.py}) orchestrates the following sequence:

\begin{enumerate}
\item \textbf{Initialization.} Load configuration files (\texttt{models.json}, \texttt{system\_prompts.json}, \texttt{test\_prompts.json}) and validate against JSON Schema. Read environment variables from \texttt{.env}. Initialize logging to console and structured JSON file (\texttt{logs/harness.log}).
\item \textbf{Capability checks.} For each enabled model, attempt to import the corresponding adapter module. If unavailable, mark the model as \texttt{dependency\_missing} and emit a warning. Perform lightweight API probes (e.g., listing available models) if supported by the provider to verify credentials and connectivity. These checks prevent silent failures mid-run.
\item \textbf{Prompt preparation.} Load prompt bank and apply optional filters (category, ID range, sampling seed). If templating is enabled, resolve Jinja2 variables using the supplied context. Record the final resolved prompts in the SQLite \texttt{prompts} table for traceability.
\item \textbf{Run metadata generation.} Create a unique \texttt{run\_id} (UUID), capture start timestamp, record Python version, dependency versions (via \texttt{pip freeze}), and price table. Store metadata in the SQLite \texttt{runs} table.
\item \textbf{Execution scheduling.} For each (model, prompt) pair, enqueue a task to the scheduler. The scheduler dispatches tasks to the worker pool respecting concurrency limits and per-provider cooldowns. Each task invokes the appropriate adapter with the resolved prompt bundle, times the call, handles retries on transient errors, and returns a standardized result dictionary.
\item \textbf{Post-processing.} For each completed call, compute derived features: response length, interrogative ratio (via regex \texttt{/\textbackslash?\$/}), keyword coverage, numeric fidelity checks. Apply the sanitization hook described in \cref{subsec:sanitization_procedures} to redact potential PII from \texttt{response\_text}. Append the enriched record to CSV and insert into SQLite.
\item \textbf{Aggregation and reporting.} After all tasks complete, compute per-model summary statistics: success rate, mean/median/p95 latency, total tokens/cost, PII flag incidence. Generate automated summary tables and plots. Emit a consolidated JSON report (\texttt{results/summary\_<run\_id>.json}) containing aggregated metrics for downstream analysis.
\item \textbf{Shutdown.} Flush output buffers, close database connections, emit final log summary, and exit with appropriate status code (0 on success, non-zero on partial completion or fatal errors).
\end{enumerate}

Error handling is structured to distinguish recoverable conditions (transient API failures, rate limits) from unrecoverable ones (missing credentials, invalid configuration). Transient errors trigger retries with exponential backoff; unrecoverable errors halt execution immediately to prevent biased results.

\subsubsection{Command-Line Interface}

The harness CLI (\texttt{python main.py [options]}) exposes flags documented in \cref{subsubsec:execution_modes}:

\begin{itemize}
\item \texttt{--models MODEL1,MODEL2}: Run only specified models (comma-separated identifiers).
\item \texttt{--category CATEGORY}: Filter prompts to a single category (e.g., \texttt{socratic\_probing}).
\item \texttt{--system-prompt TEMPLATE\_ID}: Use a specific system prompt template (default: \texttt{strict\_socratic\_v1}).
\item \texttt{--dry-run}: Validate configuration and adapters without making API calls.
\item \texttt{--seed SEED}: Set random seed for prompt sampling (default: 42).
\item \texttt{--concurrency N}: Set global worker pool size (default: 4).
\item \texttt{--output-dir DIR}: Override default output directory (\texttt{results/}).
\item \texttt{--verbose}: Enable debug-level logging for detailed diagnostics.
\end{itemize}

Example invocation for a quick check on a single model with minimal concurrency:
\begin{lstlisting}[language=bash,style=mystyle]
python main.py --models gpt-4o-mini --category socratic_probing \
  --seed 123 --concurrency 1 --verbose
\end{lstlisting}

\subsubsection{Testing and Validation}

We maintain a suite of unit and integration tests (\texttt{tests/}) covering:

\begin{itemize}
\item \textbf{Configuration validation}: Malformed JSON, missing required fields, out-of-range rate limits.
\item \textbf{Adapter contracts}: Mock API responses to verify standardized dictionary keys, error handling, cost computation.
\item \textbf{Scheduler logic}: Concurrency enforcement, cooldown timing, retry backoff sequences.
\item \textbf{Persistence schema}: SQLite table creation, foreign key constraints, CSV header alignment with schema definitions.
\item \textbf{Sanitization correctness}: PII detector recall/precision on synthetic test cases (\cref{subsec:sanitization_procedures}).
\end{itemize}

Tests are executed via \texttt{pytest} (v7.4.0) and integrated into continuous integration workflows (\cref{app:github_actions_workflow}). Code coverage is tracked via \texttt{pytest-cov} and reported in CI logs; target coverage is 85\% for core modules.

\subsubsection{Dependency Management and Reproducibility}

All dependencies are pinned in \texttt{requirements.txt} with exact versions. The harness is developed and tested on Python 3.10.12. A \texttt{Dockerfile} is provided for containerized execution, ensuring OS-level reproducibility. The container base image is \texttt{python:3.10-slim-bookworm}, and dependencies are installed from the pinned requirements file during build. Container builds are versioned and tagged to align with code releases, supporting long-term replication as described in \cref{sec:reproducibility_guide}.